\documentclass[12pt,a4paper]{article}

%% Para acentos
\usepackage[spanish,activeacute]{babel}
\usepackage[latin1]{inputenc}

% Para fuentes postscript
%\usepackage{type1ec}
%\usepackage{lmodern}
\usepackage[T1]{fontenc}


%% Para mates
\usepackage{amsmath,amssymb,latexsym}

%% Para gr�ficos
\usepackage{graphicx}

%% Para gr�ficos realizados con LatexCad
%\usepackage{latexcad}


%% Para numeraci�n autom�tica
\newcounter{nroproblema}
\setcounter{nroproblema}{1}
\newcounter{nroapartado}
\setcounter{nroapartado}{1}
\newcommand{\n}{%
    \vspace{.1in}\noindent{\bf \arabic{nroproblema})\ }%
    \addtocounter{nroproblema}{1}%
    \setcounter{nroapartado}{1}%
    }%\phantom{l}
\newcommand{\apartado}{\vspace*{0.2cm}\hspace*{0.5cm}%
   {\bf \alph{nroapartado})}
   \addtocounter{nroapartado}{1}%
   }


%% Abreviaturas comunes (para evitar usar misMacros.mcr)
\newcommand{\ds}{\displaystyle}
\newcommand{\df}{\mathop=\limits^{\sf def}}

% Para maquetaci�n
\setlength{\parindent}{0pt}
\usepackage[hmargin={1.5cm,1.5cm},vmargin={1.5cm,1.5cm}]{geometry}

\newcommand{\R}{{\mathbb{R}}}
\newcommand{\C}{{\mathbb{C}}}
\newcommand{\N}{{\mathbb{N}}}
\newcommand{\K}{{\mathbb{K}}}


\begin{document}

\begin{flushright}
1${}^o$ de grado en F�sica, curso 2013/14

\today
\end{flushright}


\begin{center}{\bf {\Large \'ALGEBRA I. HOJA 6}}\end{center}


%\begin{enumerate}


\n
Se considera el espacio vectorial $M_n(\R)$.
Calcular la dimensi\'on del subespacio
$$
    T= \{A\in M_n(\R)\colon \mathrm{tr}(A)=0\}\,.
$$

 (V\'ease el ejercicio 7 de la hoja 4 para la definici\'on de traza).










 \n  Hallar la matriz en las bases can\'onicas de la aplicaci\'on del espacio de los polinomios de grado
$\leq 3$ en $\mathbb{R}^4$ dada por $f(p(x))=(p(-1),p(1),p(-2),p(2))$.

\vspace{0.2cm}

{\it Observaci�n:} En esta hoja de ejercicios,
cuando nos referimos a vectores columna en $\mathbb{R}^n$,
omitimos el s�mbolo de transposici�n. Estrictamente hablando,
aqu\'\i \/ se trata del vector $(p(-1),p(1),p(-2),p(2))^T$.


\vspace{0.2cm}

\n Hallar las matrices en las bases can\'onicas de las siguientes aplicaciones:

\apartado La aplicaci\'on $A$ de $\mathbb{R}^3$ en $\mathbb{R}^3$ dada por
$A(x, y, z ) = (x - y + 2z , x + y, 2x + 2y + z )$.

\apartado La aplicaci\'on $B$ de $\mathbb{C}^2$ en $\mathbb{C}^2$ dada por
$B(z_1 , z_2 ) = (iz_1 + 2z_2 , z_1 - iz_2 )$.

\apartado La aplicaci\'on $C$ de $\mathbb{R}^3$ en $\mathbb{R}^3$ dada por
$C (x, y, z ) = (x + z , y, x + y + z )$.

\apartado La aplicaci\'on $D$ de $\mathbb{C}^2$ en $\mathbb{C}^3$ dada por
$D(z_1 , z_2 ) = (iz_1 + z_2 , z_1 + iz_2 , z_1- iz_2 )$.

\apartado La aplicaci\'on $E$ de $\mathbb{R}^2$ en $\mathbb{R}^4$ dada por
$E(x, y) = (x, x + y, 2x, x + 2y)$.

\vspace{0.2cm}

\n  Dada la aplicaci\'on de $\mathbb{R}^2$ en $\mathbb{R}^4$ por la matriz:
$$\left( \begin{array}{cc}
1& 0\\
0& 1\\
1& 0\\
0&1\end{array}\right)$$
en las bases can\'onicas, hallar la matriz de dicha aplicaci\'on en las bases
$\{(0, 1), (1, 1)\}$ de $\mathbb{R}^2$
y
$\{(1, 1, 0, 0), (0, 1, 1, 0), (0, 0, 1, 1)(0, 0, 0, 1)\}$ de $\mathbb{R}^4$.

\vspace{0.2cm}

\n  Dada la aplicaci\'on de $\mathbb{R}^2$ en $\mathbb{R}^4$ por la matriz:
$$\left( \begin{array}{cc}
1& 0 \\
0& 1 \\
1& 0 \\
0 &1\end{array}\right)$$
en las bases $\{(0, 1), (1, 1)\}$ de $\mathbb{R}^2$
y $\{(1, 1, 0, 0), (0, 1, 1, 0), (0, 0, 1, 1)(0, 0, 0, 1)\}$ de $\mathbb{R}^4$, hallar
la matriz de dicha aplicaci\'on en las bases can\'onicas.


\vspace{0.2cm}

\n Hallar, haciendo un cambio de base, la expresi\'on matricial en la base can\'onica de la
aplicaci\'on lineal $f$ de $\mathbb{R}^3$ en $\mathbb{R}^3$ tal que: $$f (1, 1, 0) = (1, 1, 1),\quad f (1, 0,
-1) = (0, 1, 1) \text{ y }\quad f (1, -1, 1) =
(1, 2, 2).$$


\vspace{0.2cm}

\n Determinar cu\'al es el n\'ucleo y cu\'al es la imagen:

\apartado de una proyecci\'on ortogonal de $\mathbb{R}^2$ sobre una de sus rectas.

\apartado de una proyecci\'on ortogonal de $\mathbb{R}^3$ sobre una de sus rectas.

\apartado de una proyecci\'on ortogonal de $\mathbb{R}^3$ sobre uno de sus planos.

\apartado de una simetr\'ia ortogonal de $\mathbb{R}^2$ respecto a una de sus rectas.

\apartado de una simetr\'ia ortogonal de $\mathbb{R}^3$ respecto a una de sus rectas.

\apartado de una simetr\'ia ortogonal de $\mathbb{R}^3$ respecto a uno de sus planos.


 \vspace{0.2cm}

\n Dada la aplicaci\'on lineal $f : \mathbb{R}^5 \rightarrow \mathbb{R}^4$ en las bases can\'onicas por la matriz:
$$\left( \begin{array}{ccccc}
1 &-2 &-1& 0& -1\\
1 &1 &1 &2& 0\\
2 &0 &3& 5&-3\\
0 &3& 2& 2& 1 \end{array}\right).
$$
\vskip-.2cm

\apartado Hallar una base del n\'ucleo de $f$.

\apartado Seleccionar las ecuaciones del n\'ucleo de $f$.

\apartado Hallar una base de la imagen de $f$.

\apartado Hallar las ecuaciones cartesianas de la imagen de $f$.


\vspace{0.2cm}

\n  Comprobar que es posible hallar una aplicaci\'on definida en $\mathbb{R}^3$ con imagen en $\mathbb{R}^3$, tal que:

\apartado Su n\'ucleo es el plano de ecuaci\'on $x + y + z = 0$ y su imagen es la recta $x = y = 2z$.

\apartado Su imagen es el plano de ecuaci\'on $x + y + z = 0$ y su n\'ucleo es la recta $x = y = 2z$.

\apartado Hallar las matrices correspondientes en la base can\'onica.
\vspace{0.2cm}

\n  Comprobar que es posible hallar una aplicaci\'on definida en $\mathbb{R}^4$ con imagen en $\mathbb{R}^4$, tal que
tanto su n\'ucleo como su imagen sean el subespacio de ecuaciones:
 $$\left. \begin{array}{cccccc}
x_1 +& x_2 + &x_3 +&x_4 &= 0 \\
&x_2 + &x_3 &  &= 0 \end{array}\right\}$$
Hallar la matriz de dicha aplicaci\'on en la base can\'onica de $\mathbb{R}^4$.
\vspace{0.2cm}

\n  Sea $f$ un isomorfismo en $\mathbb{R}^3$ y $g:\mathbb{R}^3\to \mathbb{R}^3$
otra aplicaci\'on lineal tal que la dimensi\'on de la imagen de
$g \circ f$ es 2.

\apartado ?`Cu\'al es el rango de la matriz de $g$?

\apartado ?`C\'ual es la dimensi\'on de la imagen de $f \circ g$?



\n Encontrar los valores de $a$ para los que los vectores
$\{(a, 0, 1), (0, 1, 1), (2, -1, a)\}$ formen una base de $\mathbb{R}^3$.

\vspace{0.2cm}





\n Comprobar que una base de $\mathbb{C}$ considerado como espacio vectorial sobre $\mathbb{C}$ est\'a formada
por $\{1\}$, pero una base de $\mathbb{C}$ considerado como espacio vectorial sobre $\mathbb{R}$ es $\{1, i\}$. Estos dos espacios
vectoriales tienen distinta dimensi\'on.

\vspace{0.2cm}

\vspace{0.2cm}



% \n Siendo
%$B_1= \{v_1=(1, 1, 0, 0), v_2=(0, 1, 0, -1), v_3=(0, 1, 1, 0), v_4=(0, 0, 0, 1)\}$
%y $B_2= \{w_1=(1, 0, 0, 1), w_2=(1, 0, 1, 0), w_3=(0, 2, 1, 0), w_4=(0, 1, 0, 1)\}$ dos bases de $\mathbb{R}^4$, hallar:
%
%\apartado Las coordenadas del vector $3w_1 + 2w_2 + w_3- w_4$ en la base $B_1=\{v_1 , v_2 , v_3 , v_4 \}$.
%
%\apartado Las coordenadas del vector $3v_1- v_3 + 2v_2$ en la base $B_2=\{w_1 , w_2 , w_3 , w_4 \}$.
%
%\apartado La matriz de cambio de base de $B_1$ a $B_2$.
%
%\apartado La matriz de cambio de base de $B_2$ a $B_1$.
%
%\vspace{0.2cm}



\n \apartado Escribir las matrices de cambio de base en el espacio de los polinomios de grado menor o igual
que $3$ con coeficientes reales entre las bases
\enspace $B_1=\{1, x-1, (x-1)^2, (x
-1)^3 \}$ \enspace y \enspace $B_2=\newline \{1, x+1, (x+1)^2, (x+1)^3\}$.

\apartado Hallar las coordenadas del polinomio $1 + x + x^2 + x^3$ directamente en las dos bases y comprobar que est\'an relacionadas por las matrices de cambio de base.





    \n Calcular  los determinantes
    siguientes:
\vskip-.3cm
    $$
a)\left|%
\begin{array}{ccc}
  0 & 2 & 4 \\
  1 & 1 & 3 \\
  3 & 3 & 3 \\
\end{array}%
\right|\ b)\left|%
\begin{array}{ccc}
  0 & -1 & 1 \\
  1 & 0 & 1 \\
  2 & -1 & -1 \\
\end{array}%
\right|\ c)\left|%
\begin{array}{ccc}
  -2 & 1 & 1 \\
  -3 & 2 & 1 \\
  1 & -1 & 1 \\
\end{array}%
\right|\ d)\left|%
\begin{array}{cccc}
  0 & -1 & -1 & 1 \\
  1 & 0 & -1 & 1 \\
  2 & -1 & -3 & -1 \\
  0 & 1 & -1 & 0 \\
\end{array}%
\right|
$$

Sol: 12, -4, -1, -8.

    \n Calcular los determinantes de las matrices dadas a
    continuaci�n:
    $$
a)\left|\frac{1}{7}\left(%
\begin{array}{ccc}
  3 & 2 & -6 \\
  6 & -3 & 2 \\
  2 & 6 & 3 \\
\end{array}%
\right) \right|
\hspace{1cm} b)\left|\frac{1}{9}\left(%
\begin{array}{ccc}
  -8 & -1 & -4 \\
  -4 & 4 & 7 \\
  1 & 8 & -4 \\
\end{array}%
\right) \right|
    $$

    Sol: a) -1,  b) 1.

%    \n Demostrar que los
%    siguientes determinantes son nulos:
%\vskip-.4cm
%    $$
%a)\left|\begin{array}{ccc}
%  7 & 1 & 3 \\
%  3 & 2 & 1 \\
%  10 & 3 & 4 \\
%\end{array}\right|\hspace{5mm}
%b)\left|\begin{array}{ccc}
%  a & b & c \\
%  -a & -b & -c \\
%  d & e & f \\
%\end{array}\right|\hspace{5mm}
%c)\left|\begin{array}{ccc}
%  1 & 2 & 3 \\
%  2 & 3 & 4 \\
%  3 & 4 & 5 \\
%\end{array}\right|\hspace{5mm}
%d)\left|\begin{array}{cccc}
%  1^2 & 2^2 & 3^2 & 4^2 \\
%  2^2 & 3^2 & 4^2 & 5^2 \\
%  3^2 & 4^2 & 5^2 & 6^2 \\
%  4^2 & 5^2 & 6^2 & 7^2 \\
%\end{array}\right|\, .
%    $$
%\vskip-.4cm


\n Sean $A,B$ dos matrices cuadradas no nulas tales que $AB=0$.

\apartado Demostrar que al menos una de las matrices $A,B$ tiene determinante nulo;

\apartado Demostrar que ambas matrices tienen determinante nulo.

    \n a) Obtener, en t�rminos de los determinantes de los
    menores diagonales de la matriz $A$,
    $$
    |A-\lambda I|=\left|\begin{array}{ccc}
      a_{11}-\lambda & a_{12} & a_{13} \\
      a_{21} & a_{22}-\lambda & a_{23} \\
      a_{31} & a_{32} & a_{33}-\lambda \\
    \end{array}
    \right|$$
    (Indicaci�n: descomponer las filas en sumandos con y sin
    $\lambda$).

    b) Aplicar la f�rmula obtenida para hallar los siguientes
    determinantes:
    $$
    |A-\lambda I|=\left|\begin{array}{ccc}
      1-\lambda & 2 & 1 \\
      0 & 3-\lambda & 1 \\
      3 & 1 & 2-\lambda \\
    \end{array}
    \right|,\hspace{1cm}\left|\frac{1}{3}\left(%
\begin{array}{ccc}
  2 & 1 & -2 \\
  -1 & -2 & -2 \\
  -2 & 2 & -1 \\
\end{array}%
\right)-\lambda I
    \right|    $$

    sol: $-{\lambda}^{3}  + {6 \lambda^{2} } - {7 \lambda} + 2$, $
-{\lambda}^{3}  - \frac{1}{3}{\lambda}^{2}  + \frac{1}{3}\lambda + 1$


    \n Sean matrices cuadradas $A\in M_{m\times m}$, $B\in M_{n\times n}$, $C\in M_{m\times
    n}$ y $D\in M_{(m+n)\times(m+n)}$ una matriz que se descompone
    por bloques como
    $$
    D=\left(%
\begin{array}{cc}
  A & C \\
  0 & B \\
\end{array}%
\right).
    $$

    a) Demostrar que $|D|=|A||B|$ utilizando la descomposici�n
    $$
    \left(%
\begin{array}{cc}
  A & C \\
  0 & B \\
\end{array}%
\right)=\left(%
\begin{array}{cc}
  I_m & 0 \\
  0 & B \\
\end{array}%
\right)\left(%
\begin{array}{cc}
  A & C \\
  0 & I_n \\
\end{array}%
\right)
    $$
    b) Demostrar, utilizando una descomposici�n similar, la misma
    igualdad $|D|=|A||B|$, en el caso
    $$
    D=\left(%
\begin{array}{cc}
  A & 0 \\
  C & B \\
\end{array}%
\right)
    $$
    (N�tese que ahora $C\in M_{n\times m}$).

    c) Calcula el determinante
$$\left| \begin{array}{cccccccc}
1 & -1 & 2 \\
   &  1  & 7 \\
   &      &  3 \\
   &      &      & 2 \\
   &      &      & 4  & -2 \\
   &      &      & 1  & 4   &  5   \\
   &      &      &     &      &      & 2 & 3 \\
   &      &      &     &      &      &    & 2
\end{array}\right|$$

Sol: -240.

%\n Estudiar si son independientes en su espacio vectorial correspondiente los siguientes conjuntos. En el caso en que sean linealmente dependientes, expresar uno de ellos como combinaci�n lineal de los restantes:
%
%\apartado $\{(1, 2, 0), (0, 1, 1)\} \subset \mathbb{R}^3$,
%
%\apartado $\{(1, 2, 0), (0, 1, 1), (1, 1, 0)\} \subset \mathbb{R}^3$,
%
%\apartado $\{(1, 2, 0), (0, 1, 1), (1, 1, -1)\} \subset \mathbb{R}^3$,
%
%\apartado $\{(1, 2, 0), (0, 1, 1), (1, 1, 0), (3, 2, 1)\} \subset \mathbb{R}^3$,
%
%\apartado $\{(i, 1 - i), (2, -2i - 2)\} \subset \mathbb{C}^2$,
%
%\apartado $\{(i, 1, -i), (1, i, -1)\} \subset \mathbb{C}^3$.
%
%\apartado \enspace $\left\{\left(\begin{matrix} 0 & 1 \\ 0 & -2\end{matrix}\right),
%\left(\begin{matrix}
%1 & -1 \\
%0 & 1
%\end{matrix}\right),
%\left(\begin{matrix}
%-2 & 5 \\
%0 & -8
%\end{matrix} \right) \right\} \subset \mathcal{M}_{2\times 2}(\mathbb{R})$
%
%
%\apartado Los polinomios $\{1, x - 2, (x - 2)^2, (x - 2)^3\} \subset P_{\mathbb{R}}^3[x]$.


\n \apartado Demostrar que si $\{f_1(x), f_2(x), \cdots , f_k (x)\}$ son funciones tales que existen $\{a_1 , a_2 ,\cdots, a_k \}$ verificando
$$\left| \begin{array}{cccc}
f_1 (a_1 ) &f_2 (a_1 ) &  \cdots & f_k (a_1 ) \\
f_1 (a_2 )& f_2 (a_2 ) & \cdots & f_k (a_2 ) \\
\vdots & \vdots & & \vdots \\
f_1(a_k) & f_2(a_k) & \cdots & f_k(a_k)
\end{array} \right| \not=0$$
entonces   $\{f_1(x), f_2(x), \cdots , f_k (x)\}$
son funciones linealmente independientes.

\vspace{0.2cm}

Demuestra los siguientes enunciados utilizando el apartado anterior:

\apartado Los polinomios
$\{1, x - 2, (x - 2)^2,\cdots, (x - 2)^n \} \subset P_{\mathbb{R}}^n[x]$
son linealmente independientes.

\apartado Para cualquier $a\in \mathbb{R}$, los polinomios
$\{1, x - a, (x - a)^2, \cdots, (x - a)^n\} \subset P_{\mathbb{R}}^n[x]$ son linealmente independientes para cualquier $n\in \mathbb{N}$.


%\n Comprobar que los siguientes conjuntos de vectores son bases de los correspondientes espacios vectoriales: (son las llamadas bases can\'onicas)

%\apartado $\{(1, 0, 0, \ldots , 0), (0, 1, 0, \ldots , 0), (0, 0, 1, \ldots , 0), \ldots , (0, 0, 0, \ldots , 1)\} \subset \mathbb{R}^n$.

%\apartado $\left\{\left(\begin{matrix}
%1 & 0\\
%0  & 0\end{matrix} \right),
%\left(\begin{matrix}
%0 & 1\\
%0 & 0\end{matrix} \right),
%\left(\begin{matrix}
%0 & 0\\
%1 & 0\end{matrix} \right),
%\left(\begin{matrix}
%0 & 0 \\
%0 & 1 \end{matrix} \right)\right\}
%\subset \mathcal{M}_{2\times2} (\mathbb{R})$

%\apartado $\{1, x, x^2,\cdots , x^n\} \subset P^n_\mathbb{R} (x)$.

\vspace{0.2cm}

\n Encontrar bases de los siguientes subespacios vectoriales de
$\mathcal{M}_{3\times3} (\mathbb{R})$:

\apartado Las matrices diagonales.

\apartado Las matrices triangulares superiores.

\apartado Las matrices triangulares inferiores.

\apartado Las matrices sim\'etricas.

\apartado Las matrices antisim\'etricas.


\n \apartado Siendo $\{e_1 , e_2 , \ldots , e_n \}$ una base de $\mathbb{R}^n$, demostrar
que los vectores:
\vspace{-.4cm}
$$
\begin{array}{cccccc}u_1 =  e_1,&
u_2 =   e_1 + e_2 ,&
\cdots, &
u_{n-1} = e_1 + e_2 + \cdots + e_{n-1} ,&
u_n =  e_1 + e_2 + \cdots + e_{n-1}+e_n
\end{array}
$$
\vskip-.3cm
son otra base de $\mathbb{R}^n$.

\apartado Hallar las coordenadas de los vectores $(n, n- 1, \ldots , 2, 1)$ y $(1, 2, \cdots , n - 1, n)$ de $\mathbb{R}^n$, en la base
$\{u_1 , u_2 , \ldots, u_n \}$, donde los $u_i$ est\'an definidos como arriba y $\{e_1 , e_2 , \ldots , e_n \}$ es la base can\'onica de $\mathbb{R}^n$.

\vspace{0.2cm}

\n Hallar las coordenadas del vector $(3, 2, -1, 0)$ de $\mathbb{R} ^4$
en la base $$B=\{(1, 0, 0, 1), (1, 0, 1, 0), (0, 2, 1, 0), (0, 1, 0, 1)\}.$$

%\vspace{0.2cm}



\vspace{0.2cm}






\end{document}
